\section{Results}
\label{sec:results}

We organize results around our four contributions. Section~\ref{sec:domain_differential} presents the flagship finding: domain-level contamination heterogeneity within MMLU that pooled analysis completely obscures. Section~\ref{sec:signal_failure} characterizes detection signal failure modes, including a directional reversal in Min-K\%++ at small model scales. Section~\ref{sec:identifiability} quantifies the structural identifiability boundaries of IRT-based contamination estimation. Section~\ref{sec:calibration_results} reports calibration quality and clarifies the value decomposition between IRT ability estimation and the contamination parameter.

\subsection{Domain-Differential Contamination}
\label{sec:domain_differential}

\paragraph{Within-benchmark heterogeneity.}
Table~\ref{tab:domain_gamma} and Figure~\ref{fig:domain_gamma} present per-domain $\gamma$ estimates from the 2PL+$\gamma$ model fitted on the $J=24$ Qwen2.5 ensemble (exp-006). MMLU exhibits striking heterogeneity across its eight cognitive domains. Math ($\gamma = +0.97$, 95\% CI $[0.53, 1.43]$, PPP$=0.44$) and Science ($\gamma = +0.84$, CI $[0.39, 1.27]$, PPP$=0.81$) show strong contamination inflation with credible intervals excluding zero and acceptable model fit. CS \& Engineering ($\gamma = +0.20$, CI $[-0.45, 0.84]$) shows a positive point estimate that does not reach significance.

The remaining five domains show null or negative $\gamma$ estimates: Business \& Law ($-0.17$), Medical ($-0.20$), Humanities ($-0.41$), and Social Sciences ($-0.50$), all with PPP$=1.00$ indicating model misfit. Negative $\gamma$ values with PPP$=1.00$ do not indicate that contamination reduces performance; rather, they signal that the 2PL+$\gamma$ specification does not capture the response pattern in these domains. One plausible explanation: these domains contain items where memorized surface patterns conflict with reasoning requirements, producing a non-monotone relationship between exposure and correctness that a single linear $\gamma$ cannot represent.

\paragraph{Dilution masking.}
The pooled MMLU estimate ($\gamma_{\text{pooled}} = +0.66$, CI $[-0.41, 1.74]$) fails to reach significance. This is not evidence of absent contamination effects. Positive STEM-domain effects are arithmetically cancelled by null and negative non-STEM estimates when aggregated, producing a near-zero pooled value. This phenomenon parallels Simpson's paradox \citep{kievit_2013_simpsons_paradox} and DIF cancellation in educational measurement \citep{ercikan_2008_dif_cancellation}: subgroup effects of opposite sign cancel at the aggregate level, creating a misleading impression of no effect. Domain-level analysis is therefore essential for contamination diagnosis; pooled estimates will systematically understate the problem when contamination is heterogeneous.

\paragraph{Cross-benchmark comparison.}
ARC-Challenge yields $\gamma = -1.46$ with a credible interval excluding zero, indicating that contaminated models perform \textit{worse} on ARC items after controlling for ability. This ``contamination forgetting'' pattern is consistent with SFT injection disrupting the reasoning strategies required for science questions. GSM8K shows $\gamma = +0.06$ (CI includes zero, PPP$=0.54$), a non-significant effect; however, GSM8K was evaluated with S1 (embedding similarity) as the only reliable signal for free-form generation items, limiting the exposure score discriminability.

\paragraph{Dose-response structure.}
Figure~\ref{fig:dose_response} visualizes the relationship between injection dosage and raw accuracy across scales. The pattern is monotonically increasing for MMLU STEM domains but shows diminishing returns above d025. At the 0.5B scale, contamination effects are catastrophic (accuracy jumps from 0.25 to 0.60 at d100 for Math items); at 3B, the same dosage produces a mild 3-5 percentage point increase. This scale-dependent attenuation is consistent with larger models having greater capacity to absorb injected items without disrupting learned representations.

% Cross-family analysis
\textit{Cross-family evidence.}
[PLACEHOLDER: pending s1 fix for full cross-family pooled IRT results.
Current v2 data (Llama-3.2-3B d000 only): ARC $\gamma=0.554$, CI$_{95}$=$[-0.121, 1.228]$; GSM8K $\gamma=-0.010$, CI$_{95}$=$[-0.441, 0.413]$.
Weak claim only: IRT calibration is robust to inclusion of a new model family, and $\gamma$ estimates remain consistent with within-family patterns.]

\subsection{Detection Signal Failure Modes}
\label{sec:signal_failure}

\paragraph{Min-K\%++ directional reversal.}
The most consequential finding for contamination detection practice is that Min-K\%++ \citep{zhang_2024_minkpp} exhibits a \textit{directional reversal} at the 0.5B scale (exp-008). Contaminated items receive \textit{lower} Min-K\%++ scores than clean items, yielding Cohen's $d = -1.55$ at d005 and $d = -1.98$ at d100 relative to d000. A practitioner applying Min-K\%++ at this scale would conclude that contaminated items are \textit{less} likely to be memorized. The signal is not merely ineffective; it is actively misleading.

At 1.5B, Min-K\%++ shows mild positive discrimination ($d \approx +0.3$); at 3B, moderate positive discrimination ($d \approx +0.7$). The transition from negative to positive discrimination occurs between 0.5B and 1.5B parameters. We hypothesize that SFT injection at 0.5B produces a distributional shift in the token-level surprisal landscape: the model's language modeling capacity is so disrupted that \textit{all} items exhibit high surprisal, but contaminated items (now overfit at the response level) paradoxically show reduced token-level fluency compared to clean items on which the model retains pretraining distributions.

\paragraph{Signal reliability across scales.}
Figure~\ref{fig:signal_heatmap} summarizes detection performance (AUC) across all signal-scale-dosage combinations. S1 (embedding similarity) achieves AUC $> 0.95$ at d005 and above for all three model scales, confirming its reliability as a standalone signal when training corpus access is available. S4 (Self-Critique) \citep{wust_2025_selfcritique} performs near-random at 0.5B (AUC $\approx 0.50$) and reaches marginal discrimination at 1.5B and 3B. RRF fusion, which assumes complementary signal contributions, degrades to S1-only performance in practice because S2 and S4 contribute noise rather than signal at small scales.

% Cross-family signal
\textit{Cross-family signal behavior.}
[PLACEHOLDER: pending s1 fix for full cross-family signal comparison. Preliminary: Llama-3.2-3B also exhibits Min-K\%++ reversal direction at comparable scales, suggesting this is a model-scale phenomenon rather than architecture-specific.]

\subsection{Identifiability Boundaries}
\label{sec:identifiability}

The simulation study (exp-004; 24 conditions, 4 AUC levels $\times$ 4 ensemble sizes $\times$ 2 contamination strengths, 10 seeds each = 240 runs) establishes the operating envelope for IRT-based contamination estimation. Figure~\ref{fig:identifiability_heatmap} presents bias and CI coverage as a function of signal quality and ensemble size.

\paragraph{Recovery thresholds.}
Parameter recovery succeeds (bias $< 0.05$, 95\% CI coverage $\geq 90\%$) at signal AUC as low as 0.55 (the minimum tested), provided $J \geq 50$ models are available. Posterior CI width decreases as $J^{-0.48}$, close to the parametric $J^{-0.5}$ rate. At $J = 24$ (our empirical ensemble size), recovery remains adequate for domains with AUC $\geq 0.70$, but CI widths increase by 40\% relative to $J = 50$.

\paragraph{Domain-method mismatch.}
Current IRT-based LLM evaluation operates at $J = 8$--$30$ models \citep{polo2026growing, lalor2024irt}. This is well below the $J \geq 50$ regime where $\gamma$ recovery is most reliable. The constraint is not a fundamental limitation of the IRT framework; it reflects the practical reality that LLM evaluation pools far fewer ``respondents'' than educational testing (where $J > 1000$ is typical). Signal quality becomes the binding constraint: when detection AUC is high ($\geq 0.90$, as S1 achieves), even $J = 16$ produces reliable $\gamma$ estimates. The identifiability boundary is therefore a joint function of $J$ and signal AUC, not either alone.

\subsection{Calibration Quality}
\label{sec:calibration_results}

\paragraph{Value decomposition.}
Table~\ref{tab:baselines} compares calibration MAE across methods. The key finding is that vanilla 2PL without a contamination parameter (B3, $\gamma = 0$) achieves calibration performance comparable to the full 2PL+$\gamma$ model across all benchmarks: ARC MAE 0.016 vs.\ 0.020, GSM8K MAE 0.006 vs.\ 0.006, MMLU MAE 0.007 vs.\ 0.007. The calibration benefit comes entirely from IRT ability estimation itself, not from the explicit $\gamma$ correction. IRT separates stable capability from item-level noise through its measurement model, producing better-calibrated $\theta$ estimates than raw accuracy regardless of whether contamination is explicitly modeled.

This result does \textit{not} undermine the ContamCal framework. The $\gamma$ parameters serve a diagnostic role distinct from calibration: they reveal \textit{which} domains exhibit inflation or deflation, the direction and magnitude of effects, and model misfit patterns. Table~\ref{tab:domain_gamma} delivers actionable information (``MMLU Math is inflated by $\sim$1 logit'') that no calibration metric can provide. Calibration answers ``how much to trust the score''; domain-level $\gamma$ answers ``where and why the score is distorted.''

\paragraph{Gating and negative results.}
ARC passes the calibration gating criteria: MAE reduction of 0.014 relative to B0-Raw (0.034 $\to$ 0.020) and a $\gamma$ estimate with CI excluding zero ($\tau = 0.756$). GSM8K fails gating because $\gamma$ is non-significant; the IRT model provides calibration improvement over raw scores but cannot claim contamination-specific correction. We also report two honest negatives. First, credible interval coverage falls below nominal levels (50\% CI achieves $\sim$35\% empirical coverage), indicating that the posterior intervals function as directional indicators rather than precise probability statements. Second, the GSM8K overcorrection ratio is 11$\times$: the model corrects 11 units of noise for every 1 unit of genuine contamination effect, a consequence of the non-significant $\gamma$ being applied uniformly.

\paragraph{Comparison to baselines.}
B0-Raw scores show substantial miscalibration on ARC (MAE$=$0.034) due to contamination-induced position bias effects, but minimal miscalibration on MMLU (MAE$=$0.016) where diverse domains partially cancel errors. B1-Remove\&Rescore reduces item count without improving calibration when contamination effects are heterogeneous across items. B4-LinearRegression achieves comparable MAE to IRT methods on GSM8K but fails on ARC where the nonlinear ability-difficulty interaction matters. B5-Emperor \citep{xu2025emperor} results confirm that mitigation-based approaches (rephrasing, option shuffling) do not restore clean performance, motivating the post-hoc calibration approach.

\begin{table}[t]
\centering
\caption{Domain-level contamination parameters from 2PL+$\gamma$ (exp-006, $J$=24 Qwen2.5 models). Domains ordered by $\gamma$ magnitude. PPP values outside $[0.1, 0.9]$ indicate model misfit.}
\label{tab:domain_gamma}
\small
\begin{tabular}{lccc}
\toprule
Domain & $\gamma$ [95\% CI] & PPP & Interpretation \\
\midrule
MMLU: Math & +0.97 [0.53, 1.43] & 0.44 & Inflation \\
MMLU: Science & +0.84 [0.39, 1.27] & 0.81 & Inflation \\
MMLU: CS/Eng & +0.20 [$-$0.45, 0.84] & 0.89 & Non-sig. \\
MMLU: Business/Law & $-$0.17 [$-$0.54, 0.23] & 1.00 & Misfit \\
MMLU: Medical & $-$0.20 [$-$0.34, $-$0.06] & 1.00 & Misfit \\
MMLU: Humanities & $-$0.41 [$-$0.78, $-$0.04] & 1.00 & Misfit \\
MMLU: Social Sci & $-$0.50 [$-$0.83, $-$0.15] & 1.00 & Misfit \\
\midrule
MMLU: Pooled & +0.66 [$-$0.41, 1.74] & n.s. & Dilution \\
ARC-Challenge & $-$1.46 [CI excl.\ 0] & -- & Forgetting \\
GSM8K & +0.06 [CI incl.\ 0] & 0.54 & Non-sig. \\
\bottomrule
\end{tabular}
\end{table}

\begin{table}[t]
\centering
\caption{Calibration comparison (MAE$\downarrow$, Kendall $\tau$$\uparrow$). B3 = vanilla 4PL ($\gamma$=0). ContamCal = 2PL+$\gamma$. Calibration benefit derives from IRT $\theta$ estimation, not from $\gamma$ correction.}
\label{tab:baselines}
\small
\begin{tabular}{lcccc}
\toprule
Method & ARC MAE & ARC $\tau$ & GSM8K MAE & MMLU MAE \\
\midrule
B0-Raw & 0.034 & -- & 0.169 & 0.016 \\
B3-Vanilla4PL & \textbf{0.016} & -- & \textbf{0.006} & \textbf{0.007} \\
ContamCal (2PL+$\gamma$) & 0.020 & 0.756 & \textbf{0.006} & \textbf{0.007} \\
\bottomrule
\end{tabular}
\end{table}
